\documentclass{article}


% if you need to pass options to natbib, use, e.g.:
%     \PassOptionsToPackage{numbers, compress}{natbib}
% before loading MURL2026


% ready for submission
\usepackage{MURL2026}


% to compile a preprint version, e.g., for submission to arXiv, add the
% [preprint] option:
%     \usepackage[preprint]{MURL2026}


% to compile a camera-ready version, add the [final] option, e.g.:
%     \usepackage[final]{MURL2026}


% to avoid loading the natbib package, add option nonatbib:
%    \usepackage[nonatbib]{MURL2026}


\usepackage[utf8]{inputenc} % allow utf-8 input
\usepackage[T1]{fontenc}    % use 8-bit T1 fonts
\usepackage{hyperref}       % hyperlinks
\usepackage{url}            % simple URL typesetting
\usepackage{booktabs}       % professional-quality tables
\usepackage{amsfonts}       % blackboard math symbols
\usepackage{nicefrac}       % compact symbols for 1/2, etc.
\usepackage{microtype}      % microtypography
\usepackage{xcolor}         % colors


\title{Formatting Instructions For MURL 2026}


% The \author macro works with any number of authors. There are two commands
% used to separate the names and addresses of multiple authors: \And and \AND.
%
% Using \And between authors leaves it to LaTeX to determine where to break the
% lines. Using \AND forces a line break at that point. So, if LaTeX puts 3 of 4
% authors names on the first line, and the last on the second line, try using
% \AND instead of \And before the third author name.


\author{%
  Author Name \\
  Department of Computer Science\\
  Maynooth University\\
  Maynooth, Co. Kildare, Ireland \\
  \texttt{author@mu.ie} \\
  % examples of more authors
  % \And
  % Coauthor \\
  % Affiliation \\
  % Address \\
  % \texttt{email} \\
  % \AND
  % Coauthor \\
  % Affiliation \\
  % Address \\
  % \texttt{email} \\
}


\begin{document}


\maketitle


\begin{abstract}
  The abstract should be a single paragraph summarising the motivation,
  methods, and key findings of your work. It must be limited to one paragraph
  and should not exceed 150 words. The word \textbf{Abstract} is centered,
  bold, and in point size 12.
\end{abstract}


\section{Submission of papers to MURL 2026}


Please read the instructions below carefully and follow them faithfully.


\subsection{Style}


Papers to be submitted to MURL 2026 must be prepared according to the
instructions presented here. Papers may only be up to {\bf six} pages long,
including figures. Additional pages \emph{containing only acknowledgments and
references} are allowed. Papers that exceed the page limit will not be
reviewed, or in any other way considered for presentation at the conference.

Authors are required to use the MURL \LaTeX{} style files provided with this
template. Please make sure you use the current files and not previous versions.
Tweaking the style files may be grounds for rejection.


\subsection{Retrieval of style files}


The style files for MURL 2026 will be made available on the conference website.
This template package contains two files that must be kept together:
\begin{itemize}
  \item \verb+MURL2026.tex+ — the main template file you write your paper in.
  \item \verb+MURL2026.sty+ — the style file that must be in the same directory
  as your \verb+.tex+ file for compilation to succeed.
\end{itemize}


\section{General formatting instructions}
\label{gen_inst}


The text must be confined within a rectangle 5.5~inches (33~picas) wide and
9~inches (54~picas) long. The left margin is 1.5~inch (9~picas). Use 10~point
type with a vertical spacing (leading) of 11~points. Times New Roman is the
preferred typeface throughout. Paragraphs are separated by \nicefrac{1}{2}~line
space (5.5 points), with no indentation.

The paper title should be 17~point, initial caps/lower case, bold, centered
between two horizontal rules. All pages should start at 1~inch (6~picas) from
the top of the page.

For the final version, authors' names are set in boldface, and each name is
centered above the corresponding address.

Please pay special attention to the instructions in Section \ref{others}
regarding figures, tables, acknowledgments, and references.


\section{Headings: first level}
\label{headings}


All headings should be lower case (except for first word and proper nouns),
flush left, and bold.

First-level headings should be in 12-point type.


\subsection{Headings: second level}


Second-level headings should be in 10-point type.


\subsubsection{Headings: third level}


Third-level headings should be in 10-point type.


\paragraph{Paragraphs}


There is also a \verb+\paragraph+ command available, which sets the heading in
bold, flush left, and inline with the text, with the heading followed by 1\,em
of space.


\section{Citations, figures, tables, references}
\label{others}


\subsection{Citations within the text}


Citations should follow standard academic practice. Use numbered citations
in square brackets, e.g.\ [1], or author-year citations where appropriate.
Ensure all references cited in the text appear in the reference list and
vice versa. We recommend using a citation management tool such as Mendeley
to keep your references organised and consistently formatted.


\subsection{Figures}


\begin{figure}
  \centering
  \fbox{\rule[-.5cm]{0cm}{4cm} \rule[-.5cm]{4cm}{0cm}}
  \caption{Sample figure caption. Figures should be self-contained: include
  in the caption all relevant parameters and hyperparameters needed to
  understand the result, so that the figure can be interpreted without
  referring back to the main text. Ensure all axis labels and legends are
  legible at the size they appear in the document.}
\end{figure}


All artwork must be neat, clean, and legible. Lines should be dark enough for
purposes of reproduction. The figure number and caption always appear after the
figure. Figures are numbered consecutively.

\textbf{Important:} Figures should be self-contained. The caption must include
all relevant parameters and hyperparameters so that the reader can fully
interpret the figure without needing to refer back to the main text. Make sure
all text, axis labels, and legends are legible at the size they appear in the
document.

You may use colour figures. However, figures should remain interpretable if
printed in black and white.


\subsection{Tables}


All tables must be centered, neat, clean and legible. The table number and
title always appear before the table. See Table~\ref{sample-table}.

Place one line space before the table title, one line space after the
table title, and one line space after the table. The table title must
be lower case (except for first word and proper nouns); tables are
numbered consecutively.

Note that publication-quality tables \emph{do not contain vertical rules.} We
strongly suggest the use of the \verb+booktabs+ package, which allows for
typesetting high-quality, professional tables.


\begin{table}
  \caption{Sample table title}
  \label{sample-table}
  \centering
  \begin{tabular}{lll}
    \toprule
    \multicolumn{2}{c}{Part}                   \\
    \cmidrule(r){1-2}
    Name     & Description     & Size ($\mu$m) \\
    \midrule
    Dendrite & Input terminal  & $\sim$100     \\
    Axon     & Output terminal & $\sim$10      \\
    Soma     & Cell body       & up to $10^6$  \\
    \bottomrule
  \end{tabular}
\end{table}


\subsection{Final instructions}


Do not change any aspects of the formatting parameters in the style files. In
particular, do not modify the width or length of the rectangle the text should
fit into, and do not change font sizes. Please note that pages should be
numbered.


\begin{ack}
Use unnumbered first level headings for the acknowledgments. All acknowledgments
go at the end of the paper before the list of references.
Do {\bf not} include this section in the anonymized submission, only in the
final paper.
\end{ack}


\section*{References}


References follow the acknowledgments. Use unnumbered first-level heading for
the references. It is permissible to reduce the font size to \verb+small+
(9 point) when listing the references. Note that the reference section does
not count towards the page limit.

\medskip

{
\small

[1] Mnih, V.\ et al.\ (2015) Human-level control through deep reinforcement
learning. {\it Nature}, 518(7540):529--533.

[2] Schulman, J.\ et al.\ (2017) Proximal policy optimization algorithms.
{\it arXiv preprint arXiv:1707.06347}.

[3] Sutton, R.S.\ \& Barto, A.G.\ (2018) {\it Reinforcement Learning: An
Introduction.} MIT Press, second edition.

}

%%%%%%%%%%%%%%%%%%%%%%%%%%%%%%%%%%%%%%%%%%%%%%%%%%%%%%%%%%%%


\end{document}
